% !PRINTED 9
\weberpage{9}
%<<CONT
mente einer stetigen geordneten Menge $\mathfrak M$, in der die
Addition besteht, und ist $c>a$, so erhält man einen Schnitt $(A,B)$
in $\mathfrak M$, wenn man alle Elemente $x$, für die $a+x\le c$ ist,
nach $A$, und für die $a+x>c$ ist, nach $B$ verweist. Dieser Schnitt
wird durch ein Element $b$ erzeugt, für das $a+b=c$ sein muss.

Die natürlichen Zahlen bilden nach unserer Definition eine messbare
Menge, in der ein kleinstes Element, nämlich $1$, vorkommt. Die
Mannigfaltigkeit der rationalen Brüche wird ebenfalls messbar, wenn
man Addition und Subtraction nach den bekannten Regeln der
Bruchrechnung erklärt. Besonders wichtig und gleichsam typisch für die
messbaren Mengen ist die Mannigfaltigkeit der geradlinigen Strecken
oder Längen einer Linie, die einfach durch Aneinanderlegen addirt
werden. Auch Stoffmengen, durch die Wage verglichen, und Zeiträume,
mit der Uhr gemessen, liefern Beispiele messbarer Mengen. Die Art des
Messens liegt nicht in der Natur der Mannigfaltigkeit selbst, sondern
wird durch den denkenden Beobachter hineingelegt; so würde es z.~B.
ebenso gut zulässig sein, unter der Summe $a+b$ zweier Strecken $a$
und $b$ die Hypotenuse eines rechtwinkligen Dreiecks mit den Katheten
$a,b$ zu verstehen, statt, wie es gewöhnlich angenommen wird, die aus
$a$ und $b$ durch Aneinanderlegen zusammengesetzte Strecke.

Um die stetige Menge $\mathfrak S$ der Schnitte in der
Mannigfaltigkeit der rationalen Brüche $\mathfrak R$ zu einer
messbaren zu machen, beachte man zunächst Folgendes.

Ist $\alpha=(A,B)$ ein Element in $\mathfrak S$ und $\mu$ ein beliebig
gegebener rationaler Bruch, so kann man immer ein Element $a'$ in $A$
so bestimmen, dass $a'+\mu=b'$ in $B$ enthalten ist. Denn wählt man
zwei beliebige Elemente $a,b$, so kann man die natürliche Zahl $m$ so
bestimmen, dass $m\mu>b-a$ ist, so dass $a+m\mu$ in $B$ enthalten ist.
Ist dann $h$ die kleinste ganze Zahl, für die $a+h\mu$ in $B$
enthalten ist, so ist $a+(h-1)\mu=a'$ in $A$ und also
$a'+\mu=b'$ in $B$ enthalten.

Wir verstehen nun unter der Summe
$(A,B)+(A',B')=(A'',B'')$ oder $\alpha+\alpha'=\alpha''$ den Schnitt
in $\mathfrak R$, den man erhält, wenn man einen rationalen Bruch
$a''$ nur dann nach $A''$ verweist, wenn ein $a$ in $A$ und ein $a'$
in $A'$ existirt von der Beschaffenheit, dass $a''\le a+a'$ ist. In
der That ist $(A'',B'')$ ein Schnitt; denn ist $a''$ in $A''$
enthalten, so gilt dasselbe von jedem kleineren Bruch, und es giebt
Brüche, die in $A''$, und Brüche, die nicht in
%CONT>>
